% TOP_PROBLEM: {"id": 2017, "title": "Erdős Problem #264", "category_id": 1, "category": {"id": 1, "name": "number_theory", "display_name": "Number Theory", "description": "Properties of integers, prime numbers, Diophantine equations.", "slug": "number-theory", "order_index": 1, "created_at": "2026-07-31T15:26:25.670Z"}, "status": "open"}
% FIRST_POSTED: null
% ENTRY_KIND: statement_only
% Imported from: unsolvedmath_additional_500.tex; UnsolvedMath ID 2017
% !TeX program = lualatex
% Compile twice: lualatex 2017.tex
% Self-contained: no external bibliography, images, or \input files.
% Numeric section labels and visible numbers are original UnsolvedMath IDs.
\documentclass[11pt,a4paper]{article}
\usepackage[margin=24mm,headheight=14pt]{geometry}
\usepackage{fontspec}
\setmainfont{Latin Modern Roman}
\setsansfont{Latin Modern Sans}
\setmonofont{Latin Modern Mono}
\usepackage{amsmath,amssymb,mathtools}
\usepackage{unicode-math}
\setmathfont{Latin Modern Math}
\usepackage{microtype}
\usepackage{enumitem,xurl,needspace}
\usepackage[dvipsnames]{xcolor}
\usepackage{hyperref}
\hypersetup{unicode=true,colorlinks=true,linkcolor=MidnightBlue,urlcolor=MidnightBlue,
 pdftitle={UnsolvedMath Problem 2017},pdfauthor={Source-annotated compilation},bookmarksnumbered=true}
\usepackage{bookmark,tocloft,titlesec,fancyhdr}
\setlength{\cftsecnumwidth}{4.2em}
\cftsetpnumwidth{2.5em}
\cftsetrmarg{4em}
\titleformat{\section}{\Large\bfseries\raggedright}{\thesection}{0.7em}{}
\pagestyle{fancy}\fancyhf{}
\fancyhead[L]{\small\sffamily Additional UnsolvedMath statements}
\fancyhead[R]{\small Original catalogue IDs}
\fancyfoot[C]{\thepage}
\setlength{\parindent}{0pt}
\setlength{\parskip}{5pt plus 1pt}
\setlength{\emergencystretch}{3em}
\setcounter{tocdepth}{1}\setcounter{secnumdepth}{1}
\setlist[itemize]{leftmargin=3em,itemsep=4pt,topsep=4pt,parsep=0pt}
\allowdisplaybreaks
\newcommand{\N}{\mathbb{N}}
\newcommand{\Z}{\mathbb{Z}}
\newcommand{\Q}{\mathbb{Q}}
\newcommand{\R}{\mathbb{R}}
\newcommand{\C}{\mathbb{C}}
\newcommand{\F}{\mathbb{F}}
\newcommand{\A}{\mathbb{A}}
\newcommand{\PP}{\mathbb{P}}
\newcommand{\E}{\mathbb{E}}
\newcommand{\eps}{\varepsilon}
\newcommand{\dd}{\,\mathrm d}
\newcommand{\sref}[2]{\hyperlink{source:#1:#2}{[S#2]}}
\newif\ifshowcatalogue
\showcataloguetrue
\newcommand{\cataloguescope}[1]{\ifshowcatalogue
 \paragraph{Statement recorded in the supplied source.}
 #1\fi}
\begin{document}
% UnsolvedMath ID 2017; source catalogue code EP-264

\setcounter{section}{2016}

\section{Irrationality under Bounded Nonzero Denominator Perturbations}\label{2017}

{\small\sffamily UnsolvedMath ID 2017 · Number Theory}

\subsection{Definitions and mathematical statement}

\paragraph{Shared notation.}Unless a section says otherwise, $\N=\{1,2,\ldots\}$,
$\Z$ is the ring of integers, $\Q,\R,\C$ are the rational, real, and complex
fields, $[N]=\{1,\ldots,\lfloor N\rfloor\}$, and $\log$ is the natural
logarithm. For real functions of a parameter tending to infinity,
$f\ll g$ and $f=O(g)$ mean $|f|\leq Cg$ eventually for some fixed $C$;
$f\gg g$ reverses this comparison; $f=o(g)$ means $f/g\to0$;
$f\sim g$ means $f/g\to1$; and $f\asymp g$ means both order bounds.
A subscript on $O$ or $\ll$ specifies allowed constant dependence.
The expression $n^{o(1)}$ in an upper bound means growth smaller than
$n^\epsilon$ for each fixed $\epsilon>0$. Local definitions govern any
exception, finite-field convention, or use of the same symbol for another
object. An asymptotic claim is not automatically an effective algorithm.



Bounded means $\sup_n|b_n|<\infty$, with the bound allowed to depend on the perturbing sequence. Every perturbation is nonzero, and no denominator may vanish. For the two proposed rapidly growing sequences, the perturbed reciprocal series converges absolutely. \sref{2017}{1} \sref{2017}{2}

\paragraph{Statement recorded in the supplied source.}
Let $a_n$ be a sequence of positive integers such that for every bounded sequence of integers $b_n$ (with $a_n+b_n\neq 0$ and $b_n\neq 0$ for all $n$) the sum $ \sum \frac{1}{a_n+b_n} $ is irrational. Are $a_n=2^n$ or $a_n=n!$ examples of such a sequence? \sref{2017}{1}

\paragraph{Residual task reported by the archive.}

The embedded review (2026-08-17) records: Determine whether the particular sequence a\_n=n! has the bounded-integer-perturbation irrationality property. \sref{2017}{1}

{\small\emph{Transcription note.} Only unambiguous escape/formatting damage and nonmathematical serialization debris were repaired in the displayed source statement.}

\subsection{Short English statement}

Can bounded nonzero integer changes to powers of two, or to factorials, ever make their reciprocal series rational?

\subsection{Sources}

{\small

\begin{itemize}

\item[\hypertarget{source:2017:1}{S1}] ULAM.AI, \emph{UnsolvedMath}, supplied \texttt{problems.json}, record 2017 (EP-264), statement and background. The embedded literature-review date is 2026-08-17; that is the archive’s review date, not a fresh certification. Dataset landing page: \url{https://huggingface.co/datasets/ulamai/UnsolvedMath}.

\item[\hypertarget{source:2017:2}{S2}] T. F. Bloom, \emph{Erdős Problems}, Problem 264, \url{https://www.erdosproblems.com/264}. Reference imported from the supplied record; the live problem page was not independently checked for this compilation.

\item[\hypertarget{source:2017:3}{S3}] V. Kovač and T. Tao, On several irrationality problems for Ahmes series, Acta Mathematica Hungarica 175 (2025), 572-608, \nolinkurl{doi:10.1007/s10474-025-01528-0.} \url{https://doi.org/10.1007/s10474-025-01528-0}. Bibliographic information imported from the supplied record.

\item[\hypertarget{source:2017:4}{S4}] Erd\"{o}s, P., On the irrationality of certain series: problems and results. New advances in transcendence theory (Durham, 1986) (1988), 102-109. Bibliographic information imported from the supplied record.

\item[\hypertarget{source:2017:5}{S5}] Erd\H{o}s, P. and Graham, R., Old and new problems and results in combinatorial number theory. Monographies de L'Enseignement Mathematique (1980). Bibliographic information imported from the supplied record.

\item[\hypertarget{source:2017:6}{S6}] Kovač, V. and Tao T., On several irrationality problems for Ahmes series. arXiv:2406.17593 (2024). Bibliographic information imported from the supplied record.

\end{itemize}

}

\label{end:2017}
\end{document}
